\chapter{Regression}\label{ch:regression}

This chapter describes the kernel regression module: the class \code{REG} for kernel
regression of a scalar function of points in $\mathbb R^d$, and the class \code{REGSURF} for
regression of (potential energy and other) surfaces of molecular geometries, including the
computation of absorption spectra by the machine-learning nuclear ensemble approach (ML-NEA).

\paragraph{Authorship and status.} The files \header{reg.h}, \header{reg.cc},
\header{regsurf.h}, \header{regsurf.cc} and the test programs \header{test_reg.cc},
\header{test_regsurf.cc} state ``Kernel ridge regression module Copyright (C) 2024 Pavel Florian
\texttt{<florian43@seznam.cz>} and Jiri Pittner'', under the same GNU GPL version~3 (or later)
licence as the rest of \LA{}. The module is written in a different style from the rest of the
library: it works mostly with raw pointers into internally allocated \code{NRVec}/\code{NRMat}
objects, returns \code{unsigned int} status codes (0 = success) in addition to calling
\code{laerror}, and reads its input from text files. The headers are \emph{not} included by
\header{la.h}; include \header{reg.h} or \header{regsurf.h} explicitly. Note that both headers
include \code{<bits/stdc++.h>} (GCC specific) and contain \code{using namespace std;} at global
scope. The classes are instantiated for \code{T=double} only.

The descriptions below say what the code does.
\lawarning{The regression module is experimental and has not yet been tested much in practice: a number of
errors were found and fixed only in October 2026. The comments in the code mention an elastic-net loss, which
the fitting algorithms do not implement; the descriptions below document what the code actually does.}

%=====================================================================
\section{Kernel functions}\label{sec:regression:kernels}
\laheader{reg.h, reg.cc}
All kernels take a difference vector \code{x} of length \code{dimensions} and a width $\sigma$.
\begin{lareference}
\lafunc{T Gaussian_kernel(const T *x, const unsigned int dimensions, const T sigma)}{$\exp(-|\mathbf x|^2/2\sigma^2)$.}
\lafunc{T Laplacian_kernel(const T *x, const unsigned int dimensions, const T sigma)}{$\exp(-\sum_i|x_i|/\sigma)$ (1-norm).}
\lafunc{T Exponential_kernel(const T *x, const unsigned int dimensions, const T sigma)}{$\exp(-|\mathbf x|/\sigma)$ (2-norm).}
\lafunc{T Matern_kernel(const T *x, const unsigned int dimensions, const T sigma, const unsigned int n, const T *Matern_comb_number)}{Mat\'ern kernel of half-integer order $\nu=n+\frac12$: $e^{-r/\sigma}\sum_{k=0}^{n} c_k\,(2r/\sigma)^{n-k}$, $r=|\mathbf x|$, i.e.\ the usual form with length scale $l=\sigma\sqrt{2\nu}$ ($n=0$ exponential, $n=1$: $(1+r/\sigma)e^{-r/\sigma}$); evaluated by the Horner scheme, $k(0)=1$. \code{Matern_comb_number} must have $n+1$ elements.}
\lafunc{unsigned int Matern_coefficient(const unsigned int n, T *comb_number)}{Fills $c_k=\frac{(n+k)!}{(2n)!}\binom nk$, $k=0..n$ ($n+1$ elements), by the recurrence $c_n=1$, $c_{k-1}=c_k\,k/((n+k)(n-k+1))$ (no factorials, no overflow for large $n$).}
\end{lareference}
The kernel is selected by name: \code{"Gaussian"}, \code{"Laplacian"}, \code{"Exponential"},
\code{"Matern"}; for Mat\'ern the order $n$ is an additional parameter (rounded to an integer).

%=====================================================================
\section{REG: kernel regression in $\mathbb R^d$}\label{sec:regression:reg}
\laheader{reg.h, reg.cc}

\subsection{Model and parameters}
The training set consists of $N$ points $\mathbf x_i\in\mathbb R^d$ (array of $N\cdot d$ values,
point after point) with values $y_i$. The model is a weighted sum of kernels centred at the
training points. The parameter array is $(\lambda, \alpha, \sigma[, n])$: regularization strength
$\lambda$, mixing $\alpha\in[0,1]$ (clamped) between the ``kernel ridge'' part $\kappa=\lambda(1-\alpha)/2$
and the ``lasso'' part $\lambda\alpha$, kernel width $\sigma$ and, for the Mat\'ern kernel, the order $n$.
As implemented, the prediction for a point $\mathbf x$ is
\[
  \hat y(\mathbf x) = \Bigl(\sum_j w_j k(\mathbf x-\mathbf x_j)\Bigr)\,\frac{S+\kappa}{S} + \lambda\alpha,
  \qquad S=\sum_j k(\mathbf x-\mathbf x_j),
\]
(the factor is omitted if $\kappa=0$ or $S=0$).

\subsection{Construction}
\begin{lareference}
\lafunc{REG(const T *x, const T *y, const T *params, const unsigned int data_size, const unsigned int x_dimensions, const string kernel)}{Model from memory (data are copied).}
\lafunc{REG(const string x_file, const string y_file, const string param_file)}{Model from text files, see below.}
\lafunc{void Init(const T *x, const T *y, const T *params, const unsigned int data_size, const unsigned int x_dimensions, const string kernel)}{Called by the constructors; allocates the working storage. With an unknown kernel name it returns silently and the object stays uninitialized (later calls fail with laerror).}
\lafunc{~REG()}{Frees the internal buffers.}
\end{lareference}
File formats (one number per line): \code{param_file} contains $N$, $d$, the kernel name,
$\lambda$, $\alpha$, $\sigma$ and for \code{"Matern"} $n$; \code{x_file} contains $N\cdot d$
coordinates, \code{y_file} $N$ values. If a file cannot be read the constructor returns silently
without initializing the object.

\subsection{Fitting and prediction}
\begin{lareference}
\lafunc{unsigned int REG_Fit(const T max_Loss, const unsigned int max_learning_cycles)}{Iterative fit. Builds the kernel matrix $K_{ij}=k(\mathbf x_i-\mathbf x_j)$ (unit diagonal) and its row sums $s_i$ on first use, starts from $w=y$ and repeats at most \code{max_learning_cycles} times the multiplicative update $p=Kw$, $p_i\leftarrow p_i(s_i+\kappa)/s_i+\lambda\alpha$, $w_i\leftarrow w_i y_i/p_i$. If \code{max_Loss}$>0$, stops when \code{Loss_Function()}$<$\code{max_Loss}.}
\lafunc{unsigned int REG_Fit_grad(const T max_Loss, const unsigned int max_learning_cycles, const T learning_rate)}{Gradient-type variant: with $p$ as above (same $\kappa$), $g_j=\sum_k(p_k-y_k)K_{jk}/\sqrt N$ and $w_j\leftarrow w_j/(1-\eta g_j)$, $\eta$=\code{learning_rate}.}
\lafunc{T REG_RMSE()}{Root mean square deviation of the predictions stored during the last fitting cycle from the training values.}
\lafunc{T Loss_Function()}{Currently identical to \code{REG_RMSE()} (the regularization terms mentioned in the comment are not added).}
\lafunc{unsigned int REG_Get_predictions(const T *x, T *y_preds, const unsigned int data_size)}{Predictions for \code{data_size} new points (array of \code{data_size}$\cdot d$ values).}
\lafunc{unsigned int REG_Get_Score(const T *x, T *y_preds, const unsigned int data_size, const T *y, T *score)}{Predictions and \code{score[0]}$=\sum_i(y_i-\hat y_i)^2$ (sum of squares, not RMSE).}
\lafunc{unsigned int REG_Save_fitted(const string model_file); unsigned int REG_Load_fitted(const string model_file)}{Write/read the $N$ weights, one per line (the model itself must be constructed with the same data and parameters).}
\lafunc{unsigned int REG_Set_params(const T *params, const string kernel); unsigned int REG_Get_params(T *params, string *kernel)}{Set/get $(\lambda,\alpha,\sigma[,n])$ and the kernel name. If the kernel, $\sigma$ or $n$ changes, the kernel matrix is rebuilt in the next fit (a change of $\lambda$, $\alpha$ only needs no rebuild).}
\lafunc{unsigned int REG_Get_weight_coeff(T *weight)}{Copy the $N$ weights.}
\lafunc{unsigned int Build_REG_distances_mat(); unsigned int Build_REG_kernels_mat(); unsigned int Build_REG_kernels_mat_sums(); T Kernel(const T *x, const unsigned int dimensions, const unsigned int kernel_id)}{Internal steps, public in the class: difference vectors, kernel matrix, its row sums, kernel by numeric id (0 Gaussian, 1 Laplacian, 2 Exponential, 3 Mat\'ern).}
\end{lareference}


\subsection{Example}
Abridged from \header{test_reg.cc}:
\begin{lstlisting}
#include "reg.h"
using namespace LA;

REG<double> reg("reg_test_x.csv", "reg_test_y.csv", "reg_test_p.csv");
reg.REG_Fit(0.0, 20);          //20 cycles, no loss threshold
double x=0.5, y;
reg.REG_Get_predictions(&x, &y, 1);
reg.REG_Save_fitted("reg_test_m.csv");
\end{lstlisting}

%=====================================================================
\section{REGSURF: regression of molecular surfaces and ML-NEA spectra}\label{sec:regression:regsurf}
\laheader{regsurf.h, regsurf.cc}

\subsection{Model}
A \code{REGSURF} model describes \code{count_levels} surfaces (e.g.\ the ground and excited state
energies) as functions of the geometry of a system of \code{count_kernels} particles (atoms) in
\code{count_dimensions} dimensions, trained on \code{count_geometries} geometries. For each
geometry the matrix of interparticle distances $r_{jk}$ is converted to a \emph{descriptor}
matrix $D_{jk}$:
\begin{itemize}
\item \code{"RE"}: $D_{jk}=r^{(0)}_{jk}/r_{jk}$, with the first training geometry as reference, zero diagonal;
\item \code{"Coulomb_matrix"}: $D_{jk}=Z_jZ_k/r_{jk}$, $D_{jj}=Z_j^{2.4}$ (needs the nuclear charges $Z$);
\item \code{"Inverse_pairs"}: $D_{jk}=1/r_{jk}$, zero diagonal;
\item a user function set by \code{s_Set_extern_descriptor}.
\end{itemize}
For geometry $g$ the kernel matrix is $K^{(g)}_{jk}=k(D^{(g)}_{jk}-D^{(0)}_{jk})$ (one-dimensional
kernel of the descriptor difference with respect to the first geometry) and its row sums
$s^{(g)}_j$ enter the prediction for level $\ell$, which is essentially $\sum_j w^{(\ell)}_j s^{(g)}_j$
with one weight per particle and level, corrected by a kernel ridge factor and the additive
$\lambda_\ell\alpha_\ell$ as in \code{REG}. The parameter array is
$(\lambda_0..\lambda_{L-1}, \alpha_0..\alpha_{L-1}, \sigma[, n])$, $L$=\code{count_levels}.

\subsection{Construction and data files}
\begin{lareference}
\lafunc{REGSURF(const unsigned int count_dimensions, const unsigned int count_kernels, const unsigned int count_geometries, const unsigned int count_levels, const string descriptor, const string kernel, const T *params, const T *geometries, const T *values, const unsigned int *Z = nullptr)}{Model from memory. \code{geometries}: $G\cdot K\cdot d$ coordinates (geometry, particle, coordinate); \code{values}: $L\cdot G$ values, level-major (all geometries of level 0 first); \code{Z}: $K$ charges (needed for \code{"Coulomb_matrix"}, otherwise \code{laerror}; may be omitted for the other descriptors).}
\lafunc{REGSURF(const string geometries_file, const string energies_file, const string parameters_file, const string Z_file = "")}{Model from text files, see below.}
\lafunc{void Init(...)}{Same arguments as the memory constructor; called by both constructors. Unknown kernel: laerror.}
\lafunc{~REGSURF()}{Frees the internal buffers.}
\end{lareference}
File formats: the parameters file contains, one per line, $d$, $K$, $G$, $L$, the descriptor
name, the kernel name, $\lambda_0..\lambda_{L-1}$, $\alpha_0..\alpha_{L-1}$, $\sigma$ and for
\code{"Matern"} $n$. The geometries file is free format: all numbers containing a decimal point
(or comma) are read in order, other tokens (atom counts, element symbols, comments) are skipped,
so concatenated XYZ files can be used. The energies file has one value per line in level-major
order, the $Z$ file one integer per line.

\subsection{Fitting, prediction and parameters}
\begin{lareference}
\lafunc{unsigned int s_Fit(const T max_Loss, const unsigned int max_learning_cycles, const T learning_rate, const bool element_invariant=true, unsigned int *kernel_classes=nullptr)}{Fits the weights of all levels: builds distances, descriptors, kernel matrices and row sums on first use, initializes the weights of level $i$ to the constant $(y_{i,0}-\lambda_i\alpha_i)/(K^2+\kappa_i)$ ($y_{i,0}$ value at the base geometry 0, $\kappa_i=\lambda_i(1-\alpha_i)K/2$) and performs at most \code{max_learning_cycles} sweeps of gradient-like updates over geometries $1..G-1$ with step \code{learning_rate}; stops a level early if \code{max_Loss}$>0$ and its loss is below \code{max_Loss}. With \code{element_invariant}, the update of each weight is replaced by the mean over the particles of the same class (\code{kernel_classes}, default the charges $Z$; without both no averaging is done).}
\lafunc{T s_RMSE(unsigned int level); T Loss_Function(unsigned int level)}{RMSE of the predictions on the training geometries for one level; \code{Loss_Function} currently returns the same.}
\lafunc{unsigned int s_Get_predictions(const unsigned int count_geometries, const unsigned int count_levels, const T *geom, const unsigned int *surfaces, T *predictions)}{Predictions for new geometries \code{geom} and the levels listed in \code{surfaces[0..count_levels-1]}; \code{predictions[i*count_levels+j]} is geometry $i$, level \code{surfaces[j]}.}
\lafunc{unsigned int s_Score(const unsigned int count_geometries, const unsigned int count_levels, const T *geom, const unsigned int *surfaces, T *predictions, const T *y_values, T *score)}{Predictions and their RMSE against \code{y_values} in \code{score[0]}.}
\lafunc{unsigned int s_Save_fitted(const string weights); unsigned int s_Load_fitted(const string weights)}{Write/read the $L\cdot K$ weights, one per line.}
\lafunc{unsigned int s_Set_params(const string descriptor, const string kernel, const T *params); unsigned int s_Get_params(string *descriptor, string *kernel, T *params)}{Change/query descriptor, kernel and parameter array; a change of descriptor, kernel, $\sigma$ or $n$ invalidates the precomputed matrices (rebuilt in the next fit; $\lambda$, $\alpha$ need no rebuild). Unknown descriptor or kernel names raise \code{laerror} (also in the constructors). The descriptor is always set from its name, so a descriptor installed by \code{s_Set_extern_descriptor} must be installed again after \code{s_Set_params}.}
\lafunc{unsigned int s_Set_extern_descriptor(unsigned int (*extern_func)(const unsigned int, const unsigned int, const T*, T*))}{Use a user descriptor function \code{f(count_particles, count_geometries, distances, desc)} which maps the $K\times K$ distance matrices of all geometries (consecutive) to descriptor matrices of the same layout.}
\lafunc{unsigned int s_Get_weight_coeff(const unsigned int s_number, T *weights)}{Copy the $K$ weights of level \code{s_number}.}
\lafunc{T Kernel(const T *x, const unsigned int dimensions, const unsigned int kernel_id); unsigned int Build_s_distances_mat(...); unsigned int Build_s_descriptor_mat(...); unsigned int Build_s_kernels_mat(...); unsigned int Build_s_kernels_mat_sums()}{Internal steps, public in the class.}
\end{lareference}
Free functions \code{Build_s_RE_descriptor_mat}, \code{Build_s_coulomb_}\allowbreak\code{matrix_descriptor_mat},
\code{Build_s_inverse_}\allowbreak\code{pairs_descriptor_mat} (\code{count_particles, count_geometries, distances, desc}
and the reference distances resp.\ $Z$) implement the built-in descriptors.

\subsection{ML-NEA absorption spectra}
\begin{lareference}
\lafunc{unsigned int s_Load_ML_NEA_geometries(const unsigned int count_geometries, const T *geometries); unsigned int s_Load_ML_NEA_geometries(const unsigned int count_geometries, const string geometries_file)}{Load an ensemble of geometries (e.g.\ sampled from a ground-state distribution) from memory or from a free-format geometry file as above.}
\lafunc{unsigned int s_Predict_ML_NEA_geometries()}{Predict all levels for all ensemble geometries.}
\lafunc{unsigned int s_Compute_ML_NEA_spectra(const T E_min, const T step, const unsigned int count_steps, const T delta)}{Spectrum on the grid $E_k=E_{\min}+k\cdot\mathrm{step}$, $k<\mathrm{count\_steps}$, as a sum over geometries and levels of Gaussians of width $\delta$ centred at the predicted excitation energies $\Delta E$ (energies assumed in atomic units), weighted by $\Delta E\,f/E$. The Gaussians are normalized, $\sqrt{2/\pi}\,\delta^{-1}\exp(-2(E-\Delta E)^2/\delta^2)$ (standard deviation $\delta/2$), and only grid points within $20\delta$ of $\Delta E$ are evaluated (beyond that the exponential underflows to zero). The oscillator strength $f=\Delta E_{\min}/\Delta E$ ($\Delta E_{\min}$ the minimum of the level over the ensemble) is only a placeholder for the true oscillator strength, which the model does not predict.}
\lafunc{unsigned int s_Save_ML_NEA_spectra(const string ML_NEA_spectra_file)}{Write the spectrum normalized to unit sum, one value per line.}
\end{lareference}

\subsection{Example}
Abridged from \header{test_regsurf.cc} (10 levels):
\begin{lstlisting}
#include "regsurf.h"
using namespace LA;

REGSURF<double> regsurf("xyz.dat", "E_sum.dat", "conf_surf.dat", "Z.dat");
regsurf.s_Fit(0.0, 100, 0.5);
for(unsigned int i=0; i<10; ++i) cout << regsurf.s_RMSE(i) << endl;
regsurf.s_Load_ML_NEA_geometries(50000, "nea_geoms.xyz");
regsurf.s_Predict_ML_NEA_geometries();
regsurf.s_Compute_ML_NEA_spectra(0.00, 0.0025, 4000, 0.2);
regsurf.s_Save_ML_NEA_spectra("spectra.dat");
\end{lstlisting}

\section{Pitfalls}\label{sec:regression:pitfalls}
\lawarning{\code{REG} and \code{REGSURF} hold raw pointers to their own buffers and have no user-defined copy constructor or assignment: copying an object leads to double deletion. Pass them by reference.}
\lawarning{Errors in the input files are partly reported by \code{laerror} and partly lead to a silently uninitialized object (\code{REG}); check the return codes and the results.}
